\section{Step 79 Average before reducing the sample}
Fill \texttt{row\_average}. Add the unsigned upper and left values and divide their sum by two with downward rounding. Put the predictor in \texttt{\$t4}. The sum can be 510, so retain it beyond eight bits until after division. A logical right shift by one gives the required floor for this nonnegative sum.

For the first sample of the worked Average row, upper is 20 and left is zero. The predictor is 10; residual 15 will restore 25. For its second red sample, left is reconstructed 25 and upper is 50, giving floor of 75 divided by two, or 37. Adding residual 18 restores 55.

The modulo-256 reduction belongs to the eventual sample store. Reducing a sum of 500 to 244 before dividing would produce 122 instead of 250. The alternative identity based on AND and XOR is valid too, provided it computes the same floor without discarding the carry.

\paragraph{Observe} Use \texttt{s79-1}; stop before \texttt{row\_average\_ready}. Inspect the state described above before continuing.

\Needspace{7\baselineskip}
\paragraph{Check} When may the Average sum be reduced to eight bits?
\begin{itemize}
  \item[A.] Before division.
  \item[B.] After the predictor has been computed and added to the residual.
  \item[C.] After either neighbor load.
\end{itemize}
